\documentclass[ekobook.tex]{subfiles}

\begin{document}\setWPP
\setlength{\multicolsep}{0em}% Remove vertical space before and after
\setlength{\columnsep}{0.5em}% Remove space between columns
\lsection{Odměňování a mzda}
\qaw{Co je mzda}{%
Mzda je \textbl{cena za práci}.
}
\qaw{Rozdíl mezi reálnou a nominální mzdou}{%
\textbl{Nominální mzda} je cifra \texthl{uvedená na výplátní pásce}, \textbl{reálná mzda} určuje, co si za ní \texthl{může zaměstnanec pořídit}.
(Je rozdíl mezi \wemph{10~000~€} a \wemph{10~000~Kč}. Je rozdíl mezi \wemph{1~000~Kč dnes} a \wemph{1~000~Kč v devatesátých letech}.)
}
\qaw{Co tvoří mzdu (základní složka, pohyblivé složky, náhrady mzdy)}{%
Mzdu tvoří:\\
\textbl{základní pevná složka}
-- smluvně daná, dle způsobu stanovení se dělí na \wemph{časovou} (hodinový plat; učitelé, úředníci),
\wemph{úkolovou} (odměna vázaná na počet výrobků),
\wemph{podílovou} (provize z výdělků; motivují se takto prodavači)
a \wemph{kombinovanou} (kombinace dříve zmíněných; číšník má daný základ + \% výdělků z tržeb)\\
\textbl{pobyhlivé složky}
-- sem patří \wemph{prémie} a \wemph{odměny}, stejně tak \wemph{osobní ohodnocení} i \wemph{příplatky ke mzdě} (za přesčasy, za rizikové podmínky, za víkendy, za svátky)\\
\textbl{náhrady mzdy}
-- sem náleží \wemph{placená dovolená} a \wemph{překážky v práci}\\
\vspace{1pt}\\
\textbl{(Nemocenská není součástí mzdy.)}
}
\qaw{Jak se počítá čistá mzda z hrubé mzdy, jaké jsou odvody státu}{%
Z hrubé mzdy se odečte \textbl{odvod zdravotního pojištění} (činí 4.5 \%),
\textbl{odvod sociálního pojištění} (činí 7.1 \%),
\textbl{daň z příjmů} (činí 15 \%), ale ta je snížena o slevy (\wemph{sleva na poplatníka}, \wemph{slevy na invaliditu/ZTP} a \wemph{sleva na děti}, kdy tato sleva jako jediná může daň snížit až do \texthl{minusu}, kdy vzniká \texthl{daňový bonus})
a ve soudně nařízených případech se můžou strhávat například alimenty.\\
\vspace{1pt}\\
\textbl{Úlevy} (jako je dárcovství krve, dary, úroky z hypotéky a další), stejně jako \wemph{sleva na manžela/ku}, se uplatňují až v ročním zúčtování v rámci \texthl{daňového přiznání}. 
}
\qaw{Co jsou slevy na dani, jak se uplatňují}{%
Až na \wemph{slevu na manžela/ku} se uplatňují \texthl{měsíčně} a snižují \texthl{daň z příjmů}.\\
\textbl{sleva na poplatníka} činí \wemph{2 570 Kč}\\
\textbl{slevy na invaliditu/ZTP} a je \wemph{210 Kč} na invaliditu I. stupně, \wemph{420 Kč} na II. stupeň a držitelé průkazu ZTP/P mají nárok na \wemph{1 345 Kč}\\
\textbl{sleva na děti} u prvního je sleva \wemph{1 267 Kč}, u druhého \wemph{1 860 Kč} a na třetí a dál \wemph{2 320 Kč};
u dětí s průkazem ZTP se sleva \texthl{zdvojnásobuje};
jako jediná může překlopit daň z příjmů do minusu a vytvoří tak \texthl{daňový bonus}\\
\textbl{sleva na manžela/ku} činí 2 070 Kč, jestli má partner ZTP, tak se sleva \texthl{zdvojnásobuje}, tato sleva se jako jediná neuplatňuje měsíčně, ale jednou ročně
}
\qaw{Co jsou úlevy na dani a kdy se uplatňují}{%
Úlevy z \texthl{daně z příjmů} se uplatňují \texthl{jednou ročně} při podání \texthl{daňového přiznání}.
Patří sem úlevy z \wemph{bezplatného dárcovství krve}, \wemph{úroky z hypotéky}, \wemph{poskytnuté dary} a \wemph{životní pojištění} i takové \wemph{důchodové připojištění}.
}
\qaw{Motivace -- hmotná, nehmotná}{%
\begin{multicols}{2}
\textbl{pozitivní hmotná}
-- prémie; služební auto; stravenky; ...\\
\textbl{pozitivní nehmotná}
-- povýšení\\
\textbl{negativní hmotná}
-- snížení mzdy; odebrání auta; ...\\
\textbl{negativní nehmotná}
-- napomenutí; sesazení z funkce; ...
\end{multicols}
}
\qaw{Výše daně z příjmů FO a PO}{%
Daňová sazba pro fyzické osoby je \textbl{15 \%} a u příjmů překračující \wemph{36násobek průměrné mzdy} (asi 1 800 000 Kč) je \textbl{23 \%}.\\
Daňová sazba pro právnické osoby je \textbl{21 \%}.
}
\qaw{Kdo podává daňové přiznání k dani z příjmu ze závislé činnosti a kdo z podnikání}{%
Ze \textbl{závislé činnosti} podávají přiznání \texthl{zaměstnanec} a \texthl{zaměstnavatel}, z \textbl{podnikání} pak \texthl{OSVČ} a to jako \wemph{propočet N x V}, \wemph{výdajovým paušálem} nebo \wemph{paušální daní}.
}
\qaw{Současný systém odvodu sociálního a zdravotního pojištění}{%
\textbl{Zdravotní pojištění} odvádíme zdravotním pojišťovnám -- ze zdravotního pojištění se platí ošetření a vyšetření u lékaře, operace, příspěvky na léky atd.\\
Pojištění platí \wemph{OSVČ}, \wemph{zaměstnanci}, \wemph{zaměstnavatelé} a \wemph{stát} (za studenty, nezaměstnané, důchodce, ženy na mateřské dovolené)\\
\vspace{1pt}\\
\textbl{Sociální pojištění} se vybírá přes \texthl{správu sociálního zabezpečení} -- vyplácí se z něj důchody a nemocenské, také se používá na \texthl{politiku nezaměstnanosti} (rekvalifikační kurzy, podpora v nezaměstnanosti...)\\
Pojištění platí \wemph{OSVČ}, \wemph{zaměstnanci} a \wemph{zaměstnavatelé}.
}
\wpage
\end{document}